\documentclass[11pt,a4paper]{article}

\usepackage[margin=1in]{geometry}
\usepackage{amsmath,amssymb,amsthm}
\usepackage{booktabs}
\usepackage{longtable}
\usepackage{array}
\usepackage{listings}
\usepackage{xcolor}
\usepackage[expansion=false]{microtype}
\usepackage{enumitem}
\usepackage[hidelinks]{hyperref}

\definecolor{codebg}{HTML}{F5F5F5}
\definecolor{warnbg}{HTML}{FFF4E5}

\lstset{
  basicstyle=\ttfamily\footnotesize,
  backgroundcolor=\color{codebg},
  frame=single,
  framesep=4pt,
  rulecolor=\color{black!20},
  breaklines=true,
  columns=fullflexible,
  keepspaces=true,
  showstringspaces=false,
  xleftmargin=0pt,
  aboveskip=1em,
  belowskip=1em
}

\newcommand{\MUST}{\textsc{must}}
\newcommand{\MUSTNOT}{\textsc{must not}}
\newcommand{\SHOULD}{\textsc{should}}
\newcommand{\SHOULDNOT}{\textsc{should not}}
\newcommand{\MAY}{\textsc{may}}

\theoremstyle{definition}
\newtheorem{definition}{Definition}
\theoremstyle{plain}
\newtheorem{proposition}{Proposition}
\theoremstyle{remark}
\newtheorem*{decision}{Decision required}

\title{\textbf{The F1R3Skein Virtual Instrument}\\[0.3em]
\large Specification, revision 2}
\author{F1R3FLY.io}
\date{4 September 2026}

\begin{document}
\maketitle

\begin{abstract}
\noindent
This document specifies the F1R3Skein virtual instrument: the gestural surface through which a
player draws melody out of the digit streams of mathematical constants, and out of which the
platform's two primary artifacts---tunes and performances---are produced. It supersedes the
behaviour described in the existing \texttt{LeapSpigotVision} README, which is inaccurate in
several respects, and it is written to be checked rather than admired: conformance language is
normative, every threshold is named and externalised, and the wire protocol is given in full.

Revision 2 makes four substantive changes to what is currently implemented. The transport moves off
the Unix domain socket, which cannot cross a device boundary. Audio moves onto the headset. The
gesture vocabulary gains explicit mutual exclusion and a calibration profile. And the semantics of
the snip gesture are corrected: as implemented, \textbf{a snipped tune is not the tune the player
heard}, and Section~\ref{sec:snip} sets out why, together with the resolution proposed for
adoption.

Section~\ref{sec:decisions} lists the seven points on which agreement is sought before
implementation begins.
\end{abstract}

\tableofcontents

\section{Scope and conformance}

\subsection{Scope}

In scope: the generative core, the instrument's state model, the gesture vocabulary and its
detection criteria, the client/engine wire protocol, the visual and audio models, the recording of
performances, and the non-functional budgets that make the instrument playable.

Out of scope, and specified elsewhere: identity and sign-on, the galleries, the engagement ledger,
breeding and selection, and the on-chain representation of tunes. This document defines what the
instrument \emph{produces}; the development plan defines what becomes of it.

\subsection{Conformance language}

\MUST, \MUSTNOT, \SHOULD, \SHOULDNOT{} and \MAY{} are used in the sense of RFC 2119. A requirement
labelled \MUST{} is a conformance condition for the September test build.

\subsection{Terminology}

\begin{description}[leftmargin=2.2cm,style=nextline,font=\normalfont\itshape]
\item[Engine] The Rust process (or on-device library) holding the authoritative instrument state.
\item[Client] The visionOS application: gesture recognition, rendering, and audio.
\item[Stream] One spigot: a constant, an output base, and a cursor position.
\item[Skein] The pair of streams, with a designation of which governs pitch and which duration.
\item[Ribbon] The visual presentation of one stream: a receding row of coloured digit patches.
\item[Tune] A snipped, named artifact. The unit of storage, breeding and gallery display.
\item[Performance] The recorded trace of a session at the instrument.
\end{description}

\section{Deployment configurations}

Two configurations are specified. Both \MUST{} be supported by the same engine code and the same
protocol; they differ only in transport and in where the engine runs.

\begin{description}[leftmargin=1cm,style=nextline]
\item[\textbf{S1 --- Tethered.}] Engine on a Mac, client on the headset, connected over TCP on the
local network with Bonjour discovery. This is the September test configuration. It is the fast path
and it is explicitly a stopgap.
\item[\textbf{S2 --- On-device.}] Engine compiled for \texttt{arm64-apple-visionos} as a static
library and called through a C foreign-function interface. No transport, no network, no discovery.
This is the target configuration.
\end{description}

The protocol of Section~\ref{sec:protocol} is defined over an abstract bidirectional message
channel. In S1 that channel is a TCP stream; in S2 it is a pair of FFI ring buffers. \textbf{No
other part of this specification distinguishes the two}, and the client \MUSTNOT{} contain
transport-specific logic outside its channel implementation.

The existing Unix-domain-socket transport \MUSTNOT{} be used for device testing. A Unix socket is
addressed by a path in one host's filesystem; the headset is a different host. It \MAY{} be
retained for engine development on the Mac.

\section{The generative core}

\subsection{Streams}

The engine \MUST{} provide the five constants already implemented---$\pi$, $e$, $\varphi$, $\ln 2$,
and Liouville's constant---each emitting digits in a configurable output base $b$ with
$2 \le b \le 36$. Digit emission \MUST{} be exact and unbounded: no fixed-precision truncation, and
no upper limit on cursor position other than time and memory.

\begin{definition}[Stream configuration]
A stream configuration is a pair $(c, b)$ of a constant and a base. Together with a cursor position
$p \in \mathbb{N}$ it determines the digit $d(c, b, p)$ uniquely.
\end{definition}

The $\pi$ implementation currently uses Gosper's unbounded linear-fractional-transformation spigot
with arbitrary-precision integers, and is materially slower than the other four. Since deep cursor
positions are exactly what a long session produces, the engine \SHOULD{} adopt a
Bailey--Borwein--Plouffe extraction for $\pi$ in bases where it applies, and \MUST{} in any case
memoise emitted digits so that re-realisation of an already-visited range is $O(1)$ per digit.

\subsection{The skein and its two cursors}

The skein holds two configurations and two \emph{independent} cursor positions, $p_L$ and $p_R$.
Independence is the whole point of the two-handed design: the left hand advances one stream, the
right hand the other, and the player's control over the music is precisely their control over the
relative rate of the two.

\begin{quote}
\textbf{Consequence, and it is load-bearing:} at any moment $p_L$ and $p_R$ are in general
\emph{different}, and the difference is the player's doing. Any operation that assumes a single
shared position is wrong. Section~\ref{sec:snip} is about the one place this assumption was made.
\end{quote}

\subsection{Note mapping}

By convention one stream governs pitch and the other duration; \emph{twist} exchanges the roles.
The engine \MUST{} produce, for the $k$-th zipped pair, a note
\[
\big(\,\text{pitch} = \Pi(d_{\text{pitch}}),\quad
      \text{duration} = \Delta(d_{\text{duration}})\,\big)
\]
where $\Pi$ is a pitch map and $\Delta$ a duration map. Both \MUST{} be parameters of the session,
not constants of the build. Pitch maps \MUST{} include at minimum major, minor, pentatonic major,
pentatonic minor, Dorian, whole tone and chromatic, each over a configurable root. Duration maps
\MUST{} include musical, linear, exponential and fixed.

The base of a stream \SHOULD{} be chosen to match the cardinality of the map it drives: a
22-element pitch vocabulary is naturally driven by a base-22 stream, a 5-element rhythmic vocabulary
by a base-5 stream. Mismatches are permitted and produce a modular fold, which is musically
legitimate but should be a choice rather than an accident.

\begin{decision}
\textbf{D-1.} Whether $\Pi$ and $\Delta$ are fixed for the duration of a performance, or are
themselves gesture-modulable mid-performance (``change scale with a gesture''). The latter is
musically far richer and is a natural extension of the vocabulary, but it adds state that the tune
representation must then carry. See Section~\ref{sec:decisions}.
\end{decision}

\section{Instrument state}

\subsection{State vector}

The engine's authoritative state \MUST{} be exactly:
\[
\Sigma = \big\langle\;
(c_L, b_L, p_L),\;
(c_R, b_R, p_R),\;
\rho,\;
\Pi,\; \Delta,\; \tau,\; \iota,\;
T
\;\big\rangle
\]
with $\rho \in \{\textsf{Stopped}, \textsf{Playing}\}$ the playback state, $\tau$ the tempo,
$\iota$ the instrument (General MIDI program), and $T$ the tune tray. The client \MUSTNOT{} hold
authoritative state; it renders $\Sigma$ and emits gestures.

\subsection{Playback state machine}

\begin{center}
\begin{tabular}{llll}
\toprule
\textbf{State} & \textbf{Event} & \textbf{Next} & \textbf{Effect} \\
\midrule
\textsf{Stopped} & \textsf{Clap}   & \textsf{Playing} & begin realisation; ribbons stitch \\
\textsf{Stopped} & \textsf{Unclap} & \textsf{Stopped} & ignored \\
\textsf{Playing} & \textsf{Unclap} & \textsf{Stopped} & halt realisation; ribbons separate \\
\textsf{Playing} & \textsf{Clap}   & \textsf{Playing} & ignored \\
any & \textsf{Pull}, \textsf{Twist}, \textsf{Scissors} & unchanged & see \S\ref{sec:gestures} \\
\bottomrule
\end{tabular}
\end{center}

Pulling \MUST{} be permitted while playing; that is how a player steers a phrase as it sounds.

\section{Gesture vocabulary}
\label{sec:gestures}

\subsection{The six gestures}

\begin{center}
\begin{tabular}{p{2.3cm} p{5.2cm} p{6.4cm}}
\toprule
\textbf{Gesture} & \textbf{Physical form} & \textbf{Effect} \\
\midrule
Pull left  & Left hand drawn toward the player & Advance the left stream; step count proportional to speed \\
Pull right & Right hand drawn toward the player & Advance the right stream, likewise \\
Twist      & One hand held clearly above the other, hands apart & Exchange the two streams, cursors included \\
Clap       & Hands brought together & Enter \textsf{Playing} \\
Unclap     & Hands drawn apart from a clap & Enter \textsf{Stopped} \\
Scissors   & Index and middle extended and spread, ring and little curled & Snip the current window into a named tune \\
\bottomrule
\end{tabular}
\end{center}

\subsection{Mutual exclusion}

The current implementation evaluates clap and twist independently, so hands brought together with
one slightly higher---an entirely natural motion---fires both. The following rules \MUST{} hold:

\begin{enumerate}[leftmargin=*]
\item Twist \MUST{} require hand separation greater than the unclap distance. Twist and clap are
      therefore never simultaneously satisfiable.
\item Scissors \MUST{} be suppressed while the hands are within clap distance of one another.
\item Pull \MUST{} be suppressed for a hand that is participating in a clap transition in the same
      frame.
\item Every gesture \MUST{} be edge-triggered with a cooldown, and \MUSTNOT{} re-fire while its
      pose is continuously held.
\end{enumerate}

\subsection{Detection parameters}

Every threshold \MUST{} live in an external calibration profile, loaded at launch and reloadable
without a rebuild. Hardcoded constants in \texttt{GestureRecognizer.swift} \MUSTNOT{} remain.

\begin{center}
\begin{tabular}{lll}
\toprule
\textbf{Parameter} & \textbf{Governs} & \textbf{Provisional} \\
\midrule
\texttt{clap.distance}        & hands-together threshold      & 0.10 m \\
\texttt{unclap.distance}      & hands-apart threshold         & 0.22 m \\
\texttt{pull.velocity\_min}   & minimum speed toward player   & 0.25 m/s \\
\texttt{pull.step\_divisor}   & speed to step count           & 0.15 \\
\texttt{pull.cooldown}        & minimum interval              & 0.08 s \\
\texttt{twist.height\_delta}  & vertical offset               & 0.06 m \\
\texttt{twist.hold\_frames}   & frames before firing          & 8 \\
\texttt{scissors.spread}      & index/middle angle            & 0.40 rad \\
\texttt{scissors.extended}    & extension floor, index/middle & \emph{to be derived} \\
\texttt{scissors.curled}      & curl ceiling, ring/little     & \emph{to be derived} \\
\texttt{scissors.hold\_frames}& frames before firing          & 5 \\
\texttt{scissors.cooldown}    & minimum interval              & 0.60 s \\
\bottomrule
\end{tabular}
\end{center}

\subsection{Finger curl must be scale-invariant}

The implemented curl metric is $\lVert \text{tip} - \text{metacarpal} \rVert / 0.08$, clamped to
unity, with ``curled'' asserted below $0.30$. On an adult hand a fully curled ring finger measures
roughly $0.06$--$0.08$ m tip-to-metacarpal, giving $0.75$--$1.0$. The threshold is unreachable, and
scissors consequently never fires from hand tracking. The metric is also not scale-invariant: the
fixed $0.08$ m divisor makes detection depend on hand size.

Curl \MUST{} instead be measured either as the interior angle at the proximal interphalangeal
joint, or as tip-to-wrist distance normalised by that finger's own summed segment lengths. Both are
scale-invariant. Thresholds \MUST{} be derived from recorded traces (\S\ref{sec:instrumentation}),
not chosen by inspection.

\section{Snip semantics}
\label{sec:snip}

\subsection{The defect}

This is the most important paragraph in the document.

The engine derives a snip as follows. \texttt{do\_snip} computes
$\textit{from} = p_L - w$ and $\textit{to} = p_L$, where $w$ is the ribbon width, and then calls
\texttt{DualStream::snip(name, from, to)}. That function creates fresh spigots from each side's
configuration, advances \emph{both} to \textit{from}, and collects $\textit{to}-\textit{from}$
pairs.

Both sides are advanced to a position computed \emph{from the left cursor alone}. But $p_L$ and
$p_R$ are independent by design---the player has been moving them at different rates all session.
So the stored snippet is the sequence
\[
\big(\, d(c_L, b_L, k),\ d(c_R, b_R, k) \,\big)_{k \,\in\, [p_L - w,\; p_L)}
\]
whereas what the player heard and saw was
\[
\big(\, d(c_L, b_L, i),\ d(c_R, b_R, j) \,\big)
\]
with $i$ and $j$ advancing independently. These coincide only if the two cursors happen to have
moved in lockstep, which is precisely the case the instrument exists to avoid.

\begin{quote}\itshape
A snipped tune is, in general, not the tune the player heard. It is a different tune, derived from
the same two constants.
\end{quote}

There is a third representation in play, which compounds the problem. The tray receives the
digits actually \emph{displayed}---\texttt{left\_ribbon.patches} zipped with
\texttt{right\_ribbon.patches}---while the snippet map receives the re-derived pairs above. So
``the snipped tune'' currently denotes three different things depending on where one looks: what
was heard, what is in the tray, and what is in the snippet map.

\subsection{Resolution}

Three resolutions are available.

\paragraph{(a) Zip-index snip.} Maintain a log of realised pairs and snip an interval of that log.
Exactly faithful to what was heard. Costs a growing log, and---decisively---destroys the address
property: a tune is no longer reconstructible from a handful of numbers, because the pairing is
arbitrary.

\paragraph{(b) Independent intervals.} \emph{Recommended.} A snip records both cursor positions and
a common length:
\[
\textsf{Snip}(k,\; i_L,\; i_R,\; n)
\quad\text{denoting}\quad
\big(\, d(c_L, b_L, i_L + t),\ d(c_R, b_R, i_R + t) \,\big)_{t \in [0, n)}.
\]
Since both ribbons are fed by the same advance operations that drive realisation, the window
$[p_L - n, p_L) \times [p_R - n, p_R)$ contains exactly the pairs displayed and sounded. Faithful,
and the address property survives: four numbers and two configurations.

\paragraph{(c) Lockstep.} Forbid independent advancement. Correct by construction and rejected: it
deletes the two-handed idea that the instrument is for.

\subsection{Consequences of (b)}

Adopting (b) requires three changes beyond the engine.

The tray, the snippet map and the realised audio \MUST{} all be derived from the same
$\textsf{Snip}(k, i_L, i_R, n)$, eliminating the three-way divergence.

The Theory of skeins requires amendment. \texttt{Skein.module} currently declares
\begin{lstlisting}
Snip . k:Skein, i:Nat, j:Nat |- "snip" "(" k "," i "," j ")" : Tune;
\end{lstlisting}
which inherits the single-interval assumption. It becomes
\begin{lstlisting}
Snip . k:Skein, iL:Nat, iR:Nat, n:Nat
     |- "snip" "(" k "," iL "," iR "," n ")" : Tune;
\end{lstlisting}
with the empty-snip equation restated as $(\textsf{Snip}\ k\ i_L\ i_R\ \textsf{NZero}) ==
(\textsf{Rest})$.

Twist then interacts correctly and for free:

\begin{proposition}
Under (b), snipping commutes with twisting: normalising $\textsf{Twist}(\textsf{Wv}\ l\ r)$ to
$\textsf{Wv}\ r\ l$ also reinterprets $i_L$ as indexing $r$, which is exactly what the engine does,
since \texttt{twist} swaps whole spigots---configuration and cursor together. No additional
equation is needed.
\end{proposition}

\begin{decision}
\textbf{D-2.} Adopt resolution (b). Confirm, and confirm the amendment to \texttt{Skein.module}.
\end{decision}

\subsection{Naming}

A scissors gesture \MUST{} suspend nothing: realisation continues, cursors keep their positions, and
the player is prompted for a name asynchronously. If no name is supplied within a timeout the
engine \MUST{} assign a provisional one rather than discard the snip. Losing a capture because a
text field was dismissed is the worst available outcome.

\section{Performance recording}

\begin{definition}[Performance trace]
An ordered sequence of timestamped records: an opening declaration of both stream configurations,
the pitch and duration maps, tempo and instrument; followed by every gesture event with its
parameters; followed by a closing timestamp.
\end{definition}

The engine \MUST{} record a trace for every session, without the player asking. Because the streams
are deterministic and the state machine is total, the trace \MUST{} replay the session exactly; no
audio is stored. Snips occurring during a session \MUST{} appear in the trace, so that a performance
resolves to the tunes harvested from it and each tune resolves back to its performance.

Trace records \MUST{} carry monotonic timestamps at millisecond resolution or better, and \SHOULD{}
additionally carry the raw hand-tracking metrics that produced each gesture, for the purpose of
\S\ref{sec:instrumentation}.

\section{Wire protocol, version 2}
\label{sec:protocol}

\subsection{Framing and encoding}

Messages \MUST{} be newline-delimited JSON, one object per line, each carrying a \texttt{type}
field. Both sides \MUST{} serialise via \texttt{Codable} and \texttt{serde} respectively; neither
side \MAY{} construct JSON by string interpolation. The current Swift client builds gesture JSON by
interpolation and escapes only the double quote, so a tune name containing a backslash produces
malformed JSON that the engine silently discards.

Every message \MUST{} carry \texttt{"v": 2}. A client receiving an unrecognised version \MUST{}
report it rather than degrade.

\subsection{Client to engine}

\begin{center}
\begin{tabular}{ll}
\toprule
\textbf{Type} & \textbf{Fields} \\
\midrule
\texttt{pull\_left}, \texttt{pull\_right} & \texttt{steps}, \texttt{velocity} \\
\texttt{twist}, \texttt{clap}, \texttt{unclap} & --- \\
\texttt{scissors} & \texttt{name} \\
\texttt{configure} & \texttt{left}, \texttt{right}, \texttt{pitch\_map}, \texttt{duration\_map}, \texttt{tempo}, \texttt{instrument} \\
\texttt{quit} & --- \\
\bottomrule
\end{tabular}
\end{center}

\texttt{configure} is new and \MUST{} be supported: the instrument is currently configured only by
command-line flags on the engine, which is unusable from a headset.

\subsection{Engine to client}

\begin{center}
\begin{tabular}{ll}
\toprule
\textbf{Type} & \textbf{Fields} \\
\midrule
\texttt{digits}    & \texttt{left}, \texttt{right}, \texttt{left\_pos}, \texttt{right\_pos} \\
\texttt{note}      & \texttt{pitch}, \texttt{duration}, \texttt{velocity}, \texttt{left\_pos}, \texttt{right\_pos}, \texttt{patch\_index} \\
\texttt{snip\_ack} & \texttt{name}, \texttt{count}, \texttt{i\_left}, \texttt{i\_right} \\
\texttt{state}     & \texttt{playing}, \texttt{left\_label}, \texttt{right\_label}, \texttt{tempo}, \texttt{instrument} \\
\texttt{status}    & \texttt{text} \\
\texttt{error}     & \texttt{code}, \texttt{text} \\
\bottomrule
\end{tabular}
\end{center}

Four changes from version 1, each fixing a defect:

\begin{enumerate}[leftmargin=*]
\item \texttt{state} replaces the practice of inferring playback state by substring-matching the
      status text. The client currently clears its playing flag on
      \texttt{statusText.contains("stopped")}, which works only until someone rewords a message.
      Playback state \MUST{} be a boolean field.
\item \texttt{note} gains \texttt{patch\_index}. The engine already computes the highlighted patch
      and does not transmit it, which is why the playing-note cursor is unimplemented on the client.
\item \texttt{snip\_ack} gains the two cursor positions, per \S\ref{sec:snip}.
\item \texttt{status} becomes advisory and human-readable only. Nothing in the client \MAY{}
      parse it.
\end{enumerate}

\subsection{Rate discipline}

The engine currently emits a \texttt{status} message on every frame of its 60\,Hz loop, beneath a
comment reading ``Push status if changed'' which describes a check the code does not perform. Each
message invalidates the entire client panel.

\begin{itemize}[leftmargin=*]
\item \texttt{status} and \texttt{state} \MUST{} be emitted only on change.
\item \texttt{digits} \MUST{} be emitted only when a cursor moves, and \MUST{} be coalesced to at
      most 30\,Hz.
\item \texttt{note} \MUST{} be emitted per note and \MUSTNOT{} be coalesced.
\item Aggregate steady-state traffic \SHOULD{} remain below 200 messages per second.
\end{itemize}

\section{Visual model}

Four elements are specified. All four exist in the current source; three are never called.

\paragraph{Ribbons.} Two receding rows of digit patches, one per stream, coloured by digit value on
a hue wheel and faded with depth. Patch entities \MUST{} be pooled and reused. Text meshes and
materials \MUST{} be cached per digit value and per depth bucket, and \MUSTNOT{} be regenerated on
every update---the current \texttt{PatchEntity.configure} calls
\texttt{MeshResource.generateText} and constructs a fresh physically-based material for every patch
on every frame, which is the dominant cost in the render path.

\paragraph{Stitching.} On entering \textsf{Playing}, threads \MUST{} animate between the two
ribbons, and \MUST{} unstitch on leaving. This is the player's only confirmation that a clap
registered, and it is currently not wired.

\paragraph{Scissor highlight.} A snip \MUST{} highlight the captured window in gold and animate its
deposit into the tray. Also implemented, also not wired.

\paragraph{Hand ghosts.} Wireframe hand skeletons \MUST{} mirror the tracked joints. The gesture
recogniser receives every hand anchor and discards it; the ghost entity is added to the scene and
never updated. Wiring these is a matter of forwarding anchors the recogniser already holds.

The tray \MUST{} display the digits actually snipped. It currently displays the live ribbon
contents, which after any subsequent pull are not the same thing.

\section{Audio model}

Audio \MUST{} be synthesised on the headset. In the current architecture the engine drives CoreMIDI
on the Mac, so the player pulls a ribbon in front of their face and hears the result from a laptop
across the room, after a network hop and a MIDI buffer. For an instrument this is disqualifying, and
it makes the latency figure---the single most important measurement of the September sessions---
unrepresentative of any shipping configuration.

\begin{itemize}[leftmargin=*]
\item The client \MUST{} render \texttt{note} messages through an on-device sampler.
\item Audio \SHOULD{} be spatialised so that each ribbon sounds from its own position in the
      immersive scene. The two streams are distinct objects in front of the player and should be
      audibly distinct.
\item Mac-side MIDI output \MAY{} be retained as an option for recording to a digital audio
      workstation, and \MUSTNOT{} be the default.
\end{itemize}

\section{Non-functional requirements}

\begin{center}
\begin{tabular}{lll}
\toprule
\textbf{Property} & \textbf{Target} & \textbf{Maximum} \\
\midrule
Gesture to audible note & 30 ms & 50 ms \\
Gesture to visual response & one frame & two frames \\
Rendered frame rate & 90 fps sustained & --- \\
Digit emission, memoised range & --- & 1 ms per 1000 digits \\
Engine to client message rate & $<$ 100 s\textsuperscript{-1} & 200 s\textsuperscript{-1} \\
Session length without degradation & 45 min & --- \\
\bottomrule
\end{tabular}
\end{center}

Failure modes \MUST{} be visible and recoverable. Loss of the channel \MUST{} surface in the panel
and \MUST{} retry with backoff, preserving engine state across reconnection. Denial of hand-tracking
authorization \MUST{} produce an explanatory message and fall back to the panel controls, rather
than a silent console line as at present. The tune tray \MUSTNOT{} silently discard entries; it
currently caps at twenty and drops the oldest, which over a long session will quietly delete the
early captures.

\section{Instrumentation}
\label{sec:instrumentation}

The gesture thresholds are guesses, one of them is provably unreachable, and there exists no data
against which to tune them. The instrument \MUST{} therefore be able to record its own inputs.

\begin{itemize}[leftmargin=*]
\item A debug overlay \MUST{} display, live, the computed metrics for each hand: separation, palm
      velocity, height differential, per-finger extension and curl, and which gesture predicates are
      currently satisfied.
\item Raw joint transforms \MUST{} be recordable to disk at tracking rate, with a manual label
      marker so a tester can annotate ``that was a scissors'' as it happens.
\item Thresholds \MUST{} be tunable offline against recorded traces and reloadable without a
      rebuild.
\end{itemize}

The first half-day of headset time \SHOULD{} be spent capturing a labelled corpus rather than
playing. Recorded traces are the durable output of the test window: they are the same artifact as a
performance trace, so this work is not a detour from the product but a first delivery of it.

\section{Conformance checklist}

A build is ready for device testing when all of the following hold.

\begin{enumerate}[leftmargin=*,itemsep=0.15em]
\item Clean install on device succeeds; entitlements and provisioning verified.
\item Hand-tracking authorization is requested, granted, and denial is handled visibly.
\item The engine is discovered automatically; the panel indicates connection within five seconds.
\item All six gestures fire, each confirmed in the debug overlay, before any music is attempted.
\item Hand ghosts are visible and track.
\item Clap produces both stitching and audio from the headset.
\item Scissors produces a name prompt; the named tune appears in the tray with the correct pair
      count and the correct digits.
\item A snipped tune, re-realised from its recorded $(k, i_L, i_R, n)$, is identical to what was
      heard at the moment of the snip.
\item 90 fps sustained with sixty patches per ribbon.
\item Tune tray and performance trace both survive an application restart.
\item Gesture traces are being written to disk.
\end{enumerate}

Item 8 is the acceptance test for Section~\ref{sec:snip} and cannot pass on the current build.

\section{Decisions sought}
\label{sec:decisions}

\begin{enumerate}[leftmargin=*,itemsep=0.4em]
\item \textbf{Snip semantics.} Adopt resolution (b), independent intervals
      $\textsf{Snip}(k, i_L, i_R, n)$, and amend \texttt{Skein.module} accordingly?
\item \textbf{Configuration ordering.} S1 tethered for September with S2 on-device to follow, or go
      straight to S2 and accept the risk to the test window? S1 is roughly three days; S2 is
      roughly a week and retires the whole transport question.
\item \textbf{Modulable maps.} Should pitch and duration maps be gesture-modulable
      mid-performance (D-1)? If yes, the tune representation must carry them, which touches the
      Theory.
\item \textbf{Ribbon width.} The snip window is currently the ribbon width, so the visible ribbon
      \emph{is} the capture buffer. Is that the intended relationship, or should the capture window
      be independently adjustable---and if so, by what gesture?
\item \textbf{Twist during playback.} Permitted, exchanging roles mid-phrase? The state machine
      above permits it. It is musically interesting and makes the resulting term slightly harder to
      read.
\item \textbf{Performance recording consent.} Recording every session by default is proposed here.
      Confirm, and confirm whether traces leave the device before the player chooses to publish.
\item \textbf{Naming timeout.} What provisional name, and after how long, when a player does not
      name a snip?
\end{enumerate}

\end{document}
